\section{Análisis dimensional de Attention}

\subsection{Dimensiones de Q, K, V}

Repasemos la fórmula básica de Attention:
$$
\text{Attention}(Q, K, V) = \text{softmax}\left(\frac{QK^T}{\sqrt{d_k}}\right) V
$$

Análisis dimensional:
\begin{itemize}
  \item $Q$: longitud de secuencia $N$, dimensión $d_k$
  \item $K, V$: longitud de secuencia $M$, dimensión $d_k$ ($K$) y $d_v$ ($V$)
  \item La longitud de secuencia de $K$ y $V$ debe coincidir (porque los pesos, tras el softmax, se multiplican por $V$)
  \item La dimensión de $Q$ debe coincidir con la de $K$ (porque hay que calcular el producto interno)
\end{itemize}

\subsection{La diferencia entre $N$ y $M$}

En la atención autorregresiva estándar $N = M$. Pero existen dos escenarios en los que $N \neq M$:
\begin{enumerate}
  \item \textbf{Atención cruzada} (Cross Attention): $Q$ proviene de una secuencia y $K, V$ de otra
  \item \textbf{Fase de Decode}: $Q$ solo tiene el token actual ($N=1$), mientras que $K, V$ contienen todos los tokens históricos ($M$ es grande)
\end{enumerate}

\subsection{Resumen de la sección}

El análisis dimensional es la base para entender la diferencia entre Prefill y Decode. La diferencia clave está en la longitud de secuencia de $Q$: en Prefill, $N$ es grande; en Decode, $N = 1$.

